\documentclass[10pt,a4paper]{amsart}
\usepackage[margin=19mm]{geometry}
\usepackage{amsmath,amssymb,mathtools}
\usepackage[scr=rsfs,cal=euler]{mathalfa}
\usepackage{xfrac}
\usepackage[unicode,hidelinks,bookmarks=false]{hyperref}
\newcommand{\ie}{i.e.\ }
\newcommand{\cf}{cf.\ }
\newcommand{\resp}{respectively\ }
\newcommand{\Subsection}{\subsection}
\providecommand{\nobreakdash}{\nobreak\hbox{-}}
\setlength{\emergencystretch}{2em}
\multlinegap0pt
\usepackage{bm}
\usepackage{bbm}
\usepackage{dsfont}
\usepackage{xspace}
\usepackage{cancel}
\usepackage{mathtools}
\usepackage{amsthm}
\makeatletter
\@ifundefined{thm}{\newtheorem{thm}{Thm}}{}
\@ifundefined{prop}{\newtheorem{prop}[thm]{Proposition}}{}
\makeatother
\title[Classifying Groups of Certain Orders]{Classifying Groups of Certain Orders}
\author{Shihan Kanungo}
\date{Source: arXiv:2606.25230v1 (CC BY 4.0)}
\begin{document}
\maketitle

\noindent\emph{Source and licence.} Classifying Groups of Certain Orders. Version arXiv:2606.25230v1 (CC BY 4.0), distributed under the Creative Commons Attribution 4.0 International licence, as recorded in the arXiv licence metadata for this exact version and shown on the abstract page https://arxiv.org/abs/2606.25230v1. Full rights notice: licenses/arxiv-2606.25230-CC-BY-4.0.txt.\par\smallskip

\noindent\textit{Editorial scope.} Counted unit: the Prop whose source label is \texttt{None} in the article (source0.tex lines 422--424, source label \texttt{None}), delivered together with the article's own complete original proof of it (source0.tex lines 425--453, one \texttt{proof} environment of 29 lines). No mathematics is written, repaired or re-derived: statement and proof are the article's own text line for line. Required context retained uncounted: none. Every definition, display and label of those lines is retained unchanged. Omitted from this package and not redistributed: 1--421, 454--603. Nothing inside a retained excerpt is omitted or altered: every retained body is delivered exactly as the article's lines are written, so no omission receipt of any kind accompanies this package. Citations: the retained excerpts carry no citation command at all, so no bibliography entry of the article is transcribed and no reference is redistributed. Reference adaptation: none was needed and none was made; every reference inside the delivered unit resolves inside it, so no unresolved reference or citation is produced. Printed slips kept literal: no printed word, symbol, hypothesis or number of the counted statement or of its proof was altered. Layout and class: the article's own class is \texttt{amsart} with packages including amsmath, amsfonts, amssymb, amsthm, amsrefs, algorithm, setspace, algpseudocode, booktabs, mathrsfs, mdframed, caption, tcolorbox, cancel; the wrapper is the lane's pinned \texttt{amsart} editorial wrapper at 10pt on A4, which restates the theorem-like environments the retained text needs and adds the article's own macro definitions that the retained text needs (none), with extra wrapper packages (bm, bbm, dsfont, xspace, cancel, mathtools). The delivered numbering is the unit's own. Assets: no retained span contains a figure environment, a graphics-inclusion command, a PDF-page inclusion, a table or a TikZ picture; no image file is copied into this package, the package declares no asset dependency and no image file is redistributed with it. Licence: version arXiv:2606.25230v1 of this article is distributed under the Creative Commons Attribution 4.0 International licence, as recorded in the arXiv licence metadata for this exact version and shown on the abstract page https://arxiv.org/abs/2606.25230v1. A full rights notice is delivered at licenses/arxiv-2606.25230-CC-BY-4.0.txt. Characters: every retained excerpt is pure ASCII, so the delivered engine, XeLaTeX with its default Latin Modern text font, reports no missing character.
\begin{prop}
Let $G$ be a group of order $p_1 p_2 \cdots p_k$, where $p_1, p_2, \dots, p_k$ are distinct primes. Then $G$ is solvable.
\end{prop}

\begin{proof}
We proceed by induction on $k$, the number of prime factors.

If $k=1$, then $|G|=p_1$ is prime, so $G$ is cyclic and hence abelian. Therefore $G$ is solvable.

Now assume $k\ge 2$ and that every group whose order is a product of fewer than $k$ distinct primes is solvable. Let $|G|=p_1p_2\cdots p_k$,
where the $p_i$ are distinct, and let $p_k$ be the largest prime dividing $|G|$. Consider the number $n_{p_k}$ of Sylow $p_k$-subgroups of $G$. By Sylow's theorems,
\[
n_{p_k}\equiv 1 \pmod{p_k} \quad \text{and} \quad 
n_{p_k}\mid p_1p_2\cdots p_{k-1}.
\]

Since $n_{p_k}$ divides the product of the primes $p_1,\dots,p_{k-1}$, every prime divisor of $n_{p_k}$ is strictly smaller than $p_k$. Thus $n_{p_k}<p_k$. 
But $n_{p_k}\equiv 1 \pmod{p_k}$, and the only positive integer less than $p_k$ that is congruent to $1$ modulo $p_k$ is $1$. Hence $n_{p_k}=1$. 
Therefore the Sylow $p_k$-subgroup $P$ is unique and hence normal in $G$.

Since $|P|=p_k$, the subgroup $P$ is cyclic of prime order and therefore abelian. In particular, $P$ is solvable. Moreover,
\[
|G/P|=p_1p_2\cdots p_{k-1},
\]
which is a product of $k-1$ distinct primes. By the induction hypothesis, the quotient group $G/P$ is solvable.
We now have a short exact sequence
\[
1 \longrightarrow P \longrightarrow G \longrightarrow G/P \longrightarrow 1,
\]
with both $P$ and $G/P$ solvable. Since an extension of a solvable group by a solvable group is solvable, it follows that $G$ is solvable.

By induction, every group of order $p_1p_2\cdots p_k$, where the $p_i$ are distinct primes, is solvable.
\end{proof}

\end{document}
